\section{The dataset: 95 recorded failures}
\label{sec:dataset}

We could not study any of this on real production outages (no ground truth), so we
built our own corpus, called \sys.

\paragraph{Two test applications, chosen to be different.} \emph{Sock Shop} is a
small online store (14 services, each with its own database). \emph{Train Ticket} is
a much larger booking system ($\sim$40 Java services) in which \emph{all} services
share one MySQL database --- a choke point that request traces cannot see inside,
which makes it the perfect stress test for the kernel layer.

\paragraph{Twelve fault types, injected on purpose.} For each recording we
deliberately broke the system in one known way, in five groups: overloading the whole
host (CPU, disk, memory, or network); capping one service's CPU or memory;
application-level faults (an artificially slowed database; a burst of connection
errors); dependency faults (freezing a service mid-flight; silently backing up a
message queue); and infrastructure faults (a ``noisy neighbour'' container stealing
CPU; degrading one service's network). Each recording stores the true answer ---
which component, what fault, when --- in a machine-readable label, plus an
independent check that the fault actually fired.

\paragraph{What one recording looks like.} Each is a four-minute session: one minute
healthy, two minutes with the fault active, one minute of recovery --- with all four
data sources recorded continuously against a single clock (aligned to about a
microsecond). In total: 95 labelled recordings, roughly 2.7\,TB raw.

\paragraph{Making kernel data usable.} The operating-system layer is recorded with
LTTng, a low-overhead Linux tracer that timestamps every system call, scheduler
decision, disk request and network event. That firehose is 2--13\,GB per four-minute
recording --- far too big to hand to any model --- so we derive a ladder of
progressively smaller views, each computed mechanically (fixed code, no AI involved)
from the level below:

\begin{itemize}
  \item \textbf{L0 --- the raw recording.} The complete binary event stream exactly
  as the tracer wrote it. It is kept archived as provenance; nothing reads it at
  analysis time except the derivation scripts below.
  \item \textbf{L1 --- per-service activity tables.} A script walks the raw stream
  event by event and attributes each event to the service that caused it, by matching
  the operating system's process identifiers against a snapshot --- captured with
  every recording --- of which container owned which processes. (This mapping step is
  what makes the data usable at all on Train Ticket, where every container identifies
  itself only as ``java''.) The attributed events are then aggregated into one row
  per service per second: typical and worst-case system-call latency, disk and
  network volume, scheduling activity, and memory-pressure signals. Gigabytes become
  a few megabytes of tables --- the ``raw numbers'' tier in our experiments.
  \item \textbf{L2 --- the wait explanation.} From the scheduler and disk events we
  reconstruct where each service's time actually went during the incident: running
  on a CPU, ready but \emph{waiting its turn} for a CPU (the signature of CPU
  starvation), waiting for disk, or blocked on something external such as a database
  or the network. These four shares answer the question activity counters cannot:
  not ``was it slow?'' but ``what was it waiting \emph{for}?''
  \item \textbf{L3 --- the plain-English digest.} A fixed template compares the
  faulty window against the healthy first minute and writes a few sentences per
  service (``system-call latency rose to 40$\times$ its baseline\ldots''), always
  with the measured numbers embedded. Because it is generated by rules rather than
  by a model, it cannot hallucinate.
\end{itemize}

The agent only ever reads L1--L3; the raw layer stays archived. Later, when we vary
``how much kernel data the agent gets,'' we are switching these rungs on and off.

\paragraph{Measured collection cost.} Recording the kernel layer costs about half a
percent of application throughput and ${\sim}13\%$ on tail latency at 200 concurrent
users --- numbers we measured ourselves rather than quoting from folklore.
